\section{Alignment}\label{Alignment}
  A creature's general moral and personal attitudes can be represented by its alignment: lawful good, neutral good, chaotic good, lawful neutral, neutral, chaotic neutral, lawful evil, neutral evil, or chaotic evil.

  Alignment is a tool for developing your character's perspective.
  Like a character's class, it is intended to provide a canvas to inspire creativity, not a narrow window to constrain identity.
  Each alignment represents a broad range of personality types or personal philosophies, so two characters of the same alignment can still be quite different from each other.
  In addition, few people are completely consistent.

  \subsection{Aligned Characters}
    Alignment is a spectrum, not a binary.
    Some characters are defined as being ``good'' or ``evil'', but this is a broad definition with a great deal of variation.
    Only angels and demons are ``pure'' good or evil.

    Approximately half of the general population of the world is neutral between good and evil, with a quarter being good and another quarter being evil.
    This does not mean that a quarter of people are all saintly do-gooders and another quarter are all psychotic murders.
    It would be more accurate to say ``good'' people are simply the most altruistic quarter of the population.
    There are a rare few saints, but all good characters have some amount of selfishness in particular contexts.
    Likewise, evil characters may act altruistically in some situations still having a fundamentally selfish nature.

    A similar ratio exists for law and chaos.
    This means that the overall population follows the ratios given in \trefnp{Aligned Population Ratios}.
    Of course, these populations are not distributed equally in the world.
    In general, humanoid civilization tends to have more good and lawful characters, and monsters tend to be more chaotic and evil.
    The GM can give more context about how alignment is used in their specific world.
    Of course, the GM can also redefine good and evil itself, so talk with them if alignment is important to you or your character.

    Monsters have their typical alignments listed in their description.
    Even monsters who are listed as ``always good'' or ``always evil'' are not \textit{maximally} good or evil -- just good or evil enough to fall into that quartile.
    Many evil monsters can be perfectly reasonable if they are capable of speech, and will uphold bargains that serve their interests.
    Likewise, good monsters have their own objectives, and will not simply do whatever an adventuring party asks of them.

    \begin{dtable}
      \lcaption{Aligned Population Ratios}
      \begin{dtabularx}{\textwidth}{X X X X}
        \tb{Alignment} & \tb{Good} & \tb{Neutral} & \tb{Evil} \tableheaderrule
        \tb{Lawful}    & 6.25\%    & 12.5\%       & 6.25\% \\
        \tb{Neutral}   & 12.5\%    & 25\%         & 12.5\% \\
        \tb{Chaotic}   & 6.25\%    & 12.5\%       & 6.25\% \\
      \end{dtabularx}
    \end{dtable}

  \subsection{Good vs. Evil}
    Distinguishing good from evil is a deeply complex task.
    In a universe where angels, demons, and deities exist and interact with the affairs of mortals, being able to clearly define good and evil is important.
    For the purposes of Rise, good and evil are strictly defined according to the intentions of one's actions, not their eventual outcomes.
    Good intentions are altruistic, and evil intentions are selfish.

    This model for good and evil has limited value as a moral system for the real world, and it neglects several dimensions of morality that people might consider important.
    It is intentionally vague about what consistutes ``other beings'', and reasonable characters might disagree about how to consider the needs of non-humanoid living things like plants and animals.
    However, is a useful way to define alignment as a character trait and roleplaying tool.
    Good characters have recognizably different behaviors from evil characters, but they are not so intrinsically opposed that they can't coexist in the same group of adventurers.
    Evil characters may cause you problems if you get in their way, but no civilized area would allow simply killing evil on sight.

    \parhead{Good} Good intentions are altruistic.
    They are based on respect and empathetic consideration for other beings.

    A good character will try to learn what other beings want or need so they can help most effectively.
    They might try to keep everyone's spirits up with cameraderie and good humor, donate money whenever possible to help those in need, or dedicate their life to punishing criminals and protecting the innocent.
    Good characters may have significant disagreements about what actions are best, and not all care about some lofty ``greater good''.
    Some believe that self-sacrifice is noble, while many would say that it does more harm than good to neglect your own needs.
    There are many interpretations of altruism.

    Even an action with good intentions may have disastrous consequences.
    Unintentionally causing harm does not make a character evil, but good characters pay attention to the effects their actions have in practice.
    If a good character caused harm, intentionally or otherwise, they would try to rectify their mistake.

    \parhead{Evil} Evil intentions are selfish.
    They come from prioritizing one's own desires when that conflicts with the knowable needs and desires of other beings.

    An evil character will generally not care what other beings want or need unless they personally benefit from that knowledge.
    They might betray allies, break laws or abuse legal loopholes to gain an advantage, or bully other creatures into doing their bidding.
    Evil characters may take actions that help others and can even work effectively as a team, but their ultimate motivation is to help themselves or make themselves feel better, not to help others.

    \parhead{Neutral} Characters that are neutral between good and evil are neither consistently altruistic nor consistently selfish.
    Most neutral characters behave altruistically in some ways and selfishly in other ways -- either at different times, or about different aspects of life.
    They often have strong bonds to particular individuals who they care about selflessly, but are not altruistic in a general sense.

    Intentions that do not involve other beings are neither good nor evil.
    Similarly, actions taken to meet one's own mandatory needs are neither good nor evil.
    Wild animals may act primarily out of self-interest, but since they generally lack the capacity to recognize the needs and desires of other beings, they are considered neutral rather than good or evil.

  \subsection{Law vs. Chaos}
    \parhead{Law} Lawful characters value consistency.
    They obey rules that guide their actions.
    Some lawful characters draw their rules from external forces, such as serving a particular master or following the legal laws of the land.
    Other lawful characters follow rules they make for themselves.

    \parhead{Chaos} Chaotic characters value flexibility and freedom.
    They make decisions based on what they think or feel at the time, even if it is inconsistent with their previous statements or actions.

    \parhead{Neutral} Characters that are neutral between law and chaos are neither exceptionally consistent nor exceptionally inconsistent.
    They tend to be generally consistent but may change their minds under the right circumstances.
    Non-sapient beings such as animals are neutral rather than lawful or chaotic.
